\appendix
\section*{Appendix C: Julia Computational Realization of Altitude Constraints}

\subsection*{C.0 Overview}

This appendix presents the Julia computational structures implementing
the altitude system of the Constellation–Manifold Toolkit.  
The purpose is to define:
\begin{itemize}
    \item altitude identifiers,
    \item allowed and forbidden transitions,
    \item purity enforcement,
    \item non-generativity enforcement,
    \item boundedness enforcement,
    \item runtime violation exceptions,
    \item the unified Series V firewall operator.
\end{itemize}

All definitions correspond directly to the mathematical semantics of
Appendix A and the Lean formalism of Appendix B.

\subsection*{C.1 Altitude Identifiers}

Julia represents altitudes using an enumeration:
\[
\texttt{@enum Altitude A0 A1 Phi7 D1 D2 D3 Phi8 Phi9 S1 I1}.
\]

This mirrors the inductive type defined in Appendix B.

\subsection*{C.2 Allowed Transition Function}

Define the computational transition predicate:
\[
T : \mathcal{A} \times \mathcal{A} \to \{0,1\}.
\]

In Julia:
\[
T(a,b) = 1 \iff (a,b) \text{ is in the allowed chain}.
\]

Explicitly:
\[
T(A_0,A_1)=1,\;
T(A_1,\Phi_7)=1,\;
T(\Phi_7,D_1)=1,\;
T(D_1,D_2)=1,\;
T(D_2,D_3)=1,
\]
\[
T(D_3,\Phi_8)=1,\;
T(\Phi_8,\Phi_9)=1,\;
T(\Phi_9,S_1)=1,\;
T(S_1,I_1)=1.
\]

Forbidden transitions:
\[
T(a,b)=0 \quad \text{otherwise}.
\]

\subsection*{C.3 Forbidden Transition Function}

Define:
\[
F(a,b) := \neg T(a,b).
\]

In Julia:
\[
\texttt{forbidden\_transition(a,b) = !allowed\_transition(a,b)}.
\]

\subsection*{C.4 Purity Enforcement}

Purity predicate:
\[
P : \mathcal{A} \to \{0,1\}.
\]

In Julia:
\[
P(a) = 1 \quad \forall a.
\]

Runtime enforcement:
\[
\texttt{enforce\_purity(a,f)} =
\begin{cases}
f(), & P(a)=1, \\
\texttt{Violation}, & P(a)=0.
\end{cases}
\]

\subsection*{C.5 Non-Generativity Enforcement}

Non-generativity predicate:
\[
N : \mathcal{A} \to \{0,1\}.
\]

In Julia:
\[
N(a)=1 \quad \forall a.
\]

Runtime enforcement:
\[
\texttt{enforce\_non\_generativity(a,f)} =
\begin{cases}
f(), & N(a)=1, \\
\texttt{Violation}, & N(a)=0.
\end{cases}
\]

\subsection*{C.6 Boundedness Enforcement}

Boundedness predicate:
\[
B : \mathcal{A} \to \{0,1\}.
\]

In Julia:
\[
B(a)=1 \quad \forall a.
\]

Runtime enforcement:
\[
\texttt{enforce\_boundedness(a,f)} =
\begin{cases}
f(), & B(a)=1, \\
\texttt{Violation}, & B(a)=0.
\end{cases}
\]

\subsection*{C.7 Runtime Violation Exception}

Define the violation type:
\[
\texttt{struct AltitudeViolation <: Exception}.
\]

This is raised whenever:
\[
F(a,b)=1 \quad \text{or} \quad P(b)=0 \quad \text{or} \quad N(b)=0 \quad \text{or} \quad B(b)=0.
\]

\subsection*{C.8 Unified Firewall Operator}

Define the computational firewall:
\[
\mathrm{Firewall}(a,b,f) =
\begin{cases}
f(), & T(a,b)=1 \land P(b) \land N(b) \land B(b), \\
\mathrm{Violation}, & \text{otherwise}.
\end{cases}
\]

In Julia:
\[
\texttt{firewall\_call(caller,callee,f)} =
\texttt{enforce\_firewall}(a,b,f).
\]

Where:
\[
a = \texttt{altitude\_of(caller)},\quad
b = \texttt{altitude\_of(callee)}.
\]

\subsection*{C.9 Altitude Mapping}

Julia maps Series IV modules to altitude identifiers:
\[
\texttt{ReversibleGeometry} \mapsto A_0,
\]
\[
\texttt{ConceptualMetric} \mapsto A_1,
\]
\[
\texttt{CurvatureEnvelope} \mapsto \Phi_7,
\]
\[
\texttt{CollapseSurface} \mapsto D_1,
\]
\[
\texttt{SoftDeformation} \mapsto D_2,
\]
\[
\texttt{EnvelopeSafety} \mapsto D_3,
\]
\[
\texttt{DriftSemigroup} \mapsto \Phi_8,
\]
\[
\texttt{InvariantExtraction} \mapsto \Phi_9,
\]
\[
\texttt{PhaseShiftOperator} \mapsto S_1,
\]
\[
\texttt{CollisionInvariant} \mapsto I_1.
\]

\subsection*{C.10 Collision Invariant Execution}

The integration altitude produces:
\[
\Omega : X \to \{0,1\}.
\]

Runtime correctness:
\[
\Omega(x)=1 \iff \text{system is collision-safe}.
\]

\[
\Omega(x)=0 \iff \text{violation detected}.
\]

This completes the computational appendix.
